\clearpage
\fancyhead[L]{\small HIND’S STAR. PARTHENOPE. VICTORIA.}
\SourceLeaf{483}
\begin{center}
\rule{0.55\linewidth}{0.4pt}\\[2pt]
Astronomische Nachrichten. Vol. XXVIII. No. 650. p. 23. 2 November 1848.\\[2pt]
\rule{0.55\linewidth}{0.4pt}
\end{center}
\begin{flushright}Göttingen, 19 August 1848.\end{flushright}

It was only at the end of July that I was able to set up the meridian circle again. I took the opportunity to observe \textsc{Hind}’s star on the meridian while it was still possible. My observation gave:

\begin{center}
1848 July 30.\quad R.A. $16^{\rm u}\ 51^{\rm m}\ 1^{\rm s}77$\\
Southern decl. $12^\circ\ 39'\ 10''6$
\end{center}

Professor \textsc{Goldschmidt} had already determined the mean position from earlier observations

\begin{center}
for 1848 Jan. 1.\quad R.A. $16^{\rm u}\ 50^{\rm m}\ 59^{\rm s}06$\\
Decl. $-12^\circ\ 39'\ 17''3$
\end{center}

and has now compared all the observations that have come to our attention with it; \mbox{---\quad---\quad---}

